\subsection*{48. Receivables and Payables: Revenue Without Cash, Work Without Payment}

Suppose an enterprise engagement has been delivered and appropriately earned ₹8 lakh of revenue but the buyer has not paid. The company carries an \textbf{accounts receivable} rather than cash. If the invoice is disputed or collection becomes doubtful, credit loss assessment and contract-specific accounting matter. Meanwhile the company must pay staff and suppliers.

Similarly, receiving supplier services before paying creates \textbf{accounts payable}. Delaying bills may temporarily preserve cash but can damage supplier trust and disguise working-capital stress. A founder should read receivable aging reports, payment terms, vendor obligations, and cash balances together. A high-revenue enterprise business with slow collections may require external working capital even when the accounting margin is attractive.

\subsection*{49. Inventory, Prepaid Expenses, and Deferred Obligations}

Physical product businesses tie cash up in raw materials, finished goods, and spare parts. Cash spent on inventory is typically not automatically the same-period cost of goods sold under accrual accounting; costs flow through inventory and are recognized when the associated goods are sold, subject to relevant measurement rules. Prepaid services and insurance may likewise create assets consumed over time.

In education and software, the mirror-image challenge is often deferred revenue or contract liabilities from customer prepayment. The cash is present but the service remains to be rendered. An annual subscription can therefore make operating cash look strong at the beginning and impose obligations long after the initial sales enthusiasm has faded.

\subsection*{50. Depreciation, Capex, and the Difference Between Expense and Investment}

Buying a computer or laboratory machine is a cash outflow. If it qualifies as a capital asset, accounting typically recognizes its cost over time through depreciation or amortization, subject to standards and useful life. This can improve reported operating income relative to immediate cash spending, but it does not undo the cash outflow. Likewise, some development costs may qualify for capitalization under specific rules while other research and development expenditure is expensed.

A founder should track \emph{capital expenditure} (capex) and the future maintenance or replacement it implies. A studio's equipment and a physical campus can tie up capital even if depreciation is gradually recognized. Never justify a capital project because annual depreciation ``looks affordable'' while ignoring purchase price, finance charges, occupancy risk, and future liquidity.

\subsection*{51. EBITDA: Useful Shorthand, Dangerous Substitute for Cash}

EBITDA refers to earnings before interest, taxes, depreciation, and amortization, commonly used as an operating proxy. It can help compare certain operating structures but is not cash flow, not statutory net profit, and not a standardized GAAP or IFRS measure. It ignores capital expenditure, working-capital demands, debt service, stock-based compensation, and other costs economically relevant to owners.

A business with high EBITDA and perpetual equipment replacement needs can consume cash. A subscription company can collect cash before earning revenue and look healthy under a different timing lens. For the early founder, \textbf{operating cash, obligations, normalized contribution, and survival runway} often matter more than a celebratory EBITDA margin shown without reconciliation.

\subsection*{52. Taxes and Indirect-Tax Collections Are Not Free Revenue}

Sales taxes such as Goods and Services Tax (GST), where applicable, involve rules about registration, classification, collection, input credit, invoicing, and remittance. Cash collected on behalf of a government is not the business's own revenue or margin. Tax obligations can create cash timing risks even in apparently profitable businesses. Corporate income tax and statutory reporting add further requirements that depend on the entity and transactions.

This chapter does not prescribe a GST classification, exemption, or Indian company-specific tax treatment for Bayes. Those require current facts and professional review. The operational principle is to separate customer fees excluding tax collected on behalf of authorities, tax liabilities, and actual spendable cash, and to avoid funding growth with amounts owed elsewhere.

\subsection*{53. The 4 Statements a Founder Should Reconcile Monthly}

Although formal reporting may require more detail, an operating founder should reconcile 4 internal statements:
\begin{enumerate}
  \item A \textbf{period profit-and-loss statement} (P\&L).
  \item A \textbf{balance sheet}.
  \item A \textbf{cash flow statement}.
  \item A \textbf{bookings, billings, and collections rollforward}.
\end{enumerate}
The 4th is a managerial bridge rather than a standardized primary financial statement. It shows how signed deals become invoices, invoices become collections, and contractual promises become recognized revenue.

If the numbers do not reconcile, investigate. Did customer cash include taxes? Did someone treat a deposit as earned revenue? Was an instructor invoice missed? Did the firm purchase an asset or incur an expense? Did a refund reduce cash and change the contract obligation? Better books do not remove uncertainty about demand, but they prevent financial uncertainty from being multiplied by inconsistent bookkeeping.
